\chapter{Inner Products and Orthogonality}
\label{ch:i13-inner-products-and-orthogonality}

Inner products add geometry to linear algebra. Length, angle, orthogonality and best approximation all come from a positive definite sesquilinear form. The complex case forces careful placement of conjugation, but the resulting theory unifies Euclidean geometry with finite-dimensional Hilbert-space methods.

\section*{Learning goals}
After this chapter, the reader should be able to:
\begin{itemize}
\item verify that a proposed formula defines an inner product;
\item compute norms, distances, and angles from an inner product;
\item prove orthogonality statements using bilinearity or sesquilinearity and conjugate symmetry;
\item compute orthogonal complements of explicitly described subspaces;
\item apply Cauchy--Schwarz and the triangle inequality with equality conditions;
\item decompose a vector relative to an orthogonal direct sum and verify the Pythagorean identity;
\end{itemize}
\section*{Conceptual roadmap}
The chapter is organized around a repeated pattern:
\[
\boxed{\text{definition}\;\longrightarrow\;\text{structural theorem}\;\longrightarrow\;\text{coordinate calculation}\;\longrightarrow\;\text{application}.}
\]
Whenever possible, we separate the invariant mathematical object from a chosen coordinate representation.

\section{Inner products}
A real inner product is symmetric bilinear and positive definite; a complex inner product is conjugate symmetric and sesquilinear.

\section{Norm and distance}
The norm is $\|v\|=\sqrt{\langle v,v\rangle}$ and the induced distance is $\|v-w\|$.

\section{Cauchy–Schwarz}
$|\langle u,v\rangle|\le\|u\|\|v\|$, with equality precisely when the vectors are linearly dependent.


% BEGIN VOL01-EXPANSION I13-example-01
\begin{example}[Cauchy--Schwarz with equality check]\label{ex:i13-expand-01}
Take \(u=(1,2,-1)\) and \(v=(2,4,-2)\). Their inner product is
\[
\langle u,v\rangle=2+8+2=12,
\]
while
\[
\|u\|=\sqrt6,\qquad \|v\|=\sqrt{24}=2\sqrt6.
\]
Hence
\[
|\langle u,v\rangle|=12=\|u\|\,\|v\|.
\]
Equality is expected because \(v=2u\). This is a useful check on both the
arithmetic and the equality condition: Cauchy--Schwarz becomes an equality
exactly along one-dimensional directions.
\end{example}
% END VOL01-EXPANSION I13-example-01
\section{Orthogonality}
Vectors are orthogonal when their inner product is zero. Orthogonal nonzero lists are automatically independent.

\section{Orthogonal complements}
For a subspace $U$, $U^\perp=\{v:\langle v,u\rangle=0\ \forall u\in U\}$. In finite dimensions, $V=U\oplus U^\perp$.


% BEGIN VOL01-EXPANSION I13-example-02
\begin{example}[Computing an orthogonal complement]\label{ex:i13-expand-02}
Let
\[
U=\operatorname{span}\{(1,1,0),(0,1,1)\}\subset\mathbb R^3.
\]
A vector \(x=(a,b,c)\) lies in \(U^\perp\) exactly when
\[
a+b=0,\qquad b+c=0.
\]
Thus \(a=-b\) and \(c=-b\), so
\[
U^\perp=\operatorname{span}\{(-1,1,-1)\}.
\]
The dimension check confirms the result:
\(\dim U=2\), \(\dim U^\perp=1\), and the dimensions add to \(3\).
\end{example}
% END VOL01-EXPANSION I13-example-02
\section{Parallelogram and polarization}
The inner product determines the norm, and conversely a norm comes from an inner product exactly when it satisfies the parallelogram law.


% BEGIN VOL01-EXPANSION I13-example-03
\begin{example}[Polarization recovers hidden geometry]\label{ex:i13-expand-03}
Suppose a real inner-product norm is known but the inner product is not.
For
\[
u=(2,1),\qquad v=(1,-1),
\]
the polarization identity gives
\[
\langle u,v\rangle
=\frac14\bigl(\|u+v\|^2-\|u-v\|^2\bigr).
\]
Here \(u+v=(3,0)\) and \(u-v=(1,2)\), so
\[
\langle u,v\rangle=\frac14(9-5)=1.
\]
Direct computation gives \(2\cdot1+1(-1)=1\), providing an independent check.
The norm therefore encodes the inner-product geometry when the parallelogram
law holds.
\end{example}
% END VOL01-EXPANSION I13-example-03
\section{Core structural results}

\begin{theorem}[Cauchy–Schwarz inequality]\label{thm:i13-01}
For all $u,v$ in an inner-product space, $|\langle u,v\rangle|\le\|u\|\|v\|$.
\begin{proof}
If $v\ne0$, apply positivity to $u-(\langle u,v\rangle/\|v\|^2)v$ and expand. The case $v=0$ is immediate.
\end{proof}
\end{theorem}

\begin{theorem}[Pythagorean theorem]\label{thm:i13-02}
If $u\perp v$, then $\|u+v\|^2=\|u\|^2+\|v\|^2$.
\begin{proof}
Expand $\langle u+v,u+v\rangle$ and use conjugate symmetry together with the vanishing cross terms.
\end{proof}
\end{theorem}

\section{Worked examples}

\begin{example}[Euclidean product]\label{ex:i13-worked-01}
On $\mathbb R^n$, $\langle x,y\rangle=x^Ty$.
\end{example}

\begin{example}[Weighted product]\label{ex:i13-worked-02}
For positive weights $w_i$, $\langle x,y\rangle=\sum_i w_ix_iy_i$ is an inner product.
\end{example}

\begin{example}[Function product]\label{ex:i13-worked-03}
On continuous functions, $\langle f,g\rangle=\int_0^1 f(t)\overline{g(t)}\,dt$ is an inner product.
\end{example}

\section{Diagnostic checklist}
Before moving on, it is useful to distinguish four levels of understanding:
\begin{enumerate}
\item recognize the definitions in unfamiliar examples;
\item prove the core structural statements without coordinates when possible;
\item translate the same statement into a matrix or coordinate computation;
\item know which hypotheses are essential and produce a counterexample when one is removed.
\end{enumerate}

% BEGIN VOLUME01 SOLVED DOSSIERS
\section{Solved dossiers}
These dossiers are longer worked problems. They complement the shorter exercise--hint--solution triads by combining several ideas in one argument.

\begin{problem}[Cauchy–Schwarz equality]\label{prob:i13-dossier-01}
When does equality hold in Cauchy–Schwarz?
\end{problem}
\begin{solution}
Equality holds exactly when one vector is a scalar multiple of the other. In the standard proof, equality means the orthogonal residual $u-(\langle u,v\rangle/\|v\|^2)v$ has norm zero.
\end{solution}

\begin{problem}[Orthogonal independence]\label{prob:i13-dossier-02}
Prove nonzero pairwise orthogonal vectors are linearly independent.
\end{problem}
\begin{solution}
Take $\sum a_iv_i=0$ and inner-product with $v_j$. All cross terms vanish, leaving $a_j\|v_j\|^2=0$. Since $v_j\ne0$, $a_j=0$ for every $j$.
\end{solution}

\begin{problem}[Weighted inner product]\label{prob:i13-dossier-03}
Show $\langle x,y\rangle=2x_1y_1+5x_2y_2$ is an inner product on $\mathbb R^2$.
\end{problem}
\begin{solution}
Bilinearity and symmetry are immediate. Positive definiteness follows because $\langle x,x\rangle=2x_1^2+5x_2^2>0$ for every nonzero $x$.
\end{solution}

\begin{problem}[Orthogonal complement]\label{prob:i13-dossier-04}
Find the orthogonal complement of $\operatorname{span}(1,1,1)$ in $\mathbb R^3$.
\end{problem}
\begin{solution}
It is the plane $\{(x,y,z):x+y+z=0\}$, since orthogonality to $(1,1,1)$ is exactly the displayed equation.
\end{solution}

\begin{problem}[Pythagoras]\label{prob:i13-dossier-05}
If $u\perp v$, prove $\|u+v\|^2=\|u\|^2+\|v\|^2$.
\end{problem}
\begin{solution}
Expand $\langle u+v,u+v\rangle$. The cross terms vanish by orthogonality, leaving the sum of the squared norms.
\end{solution}

\begin{problem}[Triangle inequality]\label{prob:i13-dossier-06}
Derive the triangle inequality from Cauchy–Schwarz.
\end{problem}
\begin{solution}
$\|u+v\|^2=\|u\|^2+2\operatorname{Re}\langle u,v\rangle+\|v\|^2\le(\|u\|+\|v\|)^2$. Taking nonnegative square roots yields the result.
\end{solution}

\begin{problem}[Complex conjugation placement]\label{prob:i13-dossier-07}
Why must the standard complex inner product include conjugation?
\end{problem}
\begin{solution}
Without conjugation, $\sum z_i^2$ can be nonreal or negative-looking; for example $i^2=-1$. With conjugation, $\langle z,z\rangle=\sum|z_i|^2\ge0$ and is zero only for $z=0$.
\end{solution}

\begin{problem}[Distance to a subspace]\label{prob:i13-dossier-08}
If $V=U\oplus U^\perp$ and $v=u+w$, which component controls the distance from $v$ to $U$?
\end{problem}
\begin{solution}
For any $u'\in U$, $v-u'=(u-u')+w$ with orthogonal terms, so $\|v-u'\|^2=\|u-u'\|^2+\|w\|^2$. The minimal distance is $\|w\|$.
\end{solution}

\begin{problem}[Polarization]\label{prob:i13-dossier-09}
Recover a real inner product from its norm.
\end{problem}
\begin{solution}
The real polarization identity is $\langle u,v\rangle=\frac14(\|u+v\|^2-\|u-v\|^2)$. Expanding the two norms verifies the formula.
\end{solution}

\begin{problem}[Angle computation]\label{prob:i13-dossier-10}
Find the angle between $(1,1,0)$ and $(1,0,1)$.
\end{problem}
\begin{solution}
The dot product is $1$, each norm is $\sqrt2$, so $\cos\theta=1/2$ and $\theta=\pi/3$.
\end{solution}

\begin{problem}[Orthogonal decomposition dimension]\label{prob:i13-dossier-11}
If $U$ is a subspace of finite-dimensional $V$, prove $\dim U+\dim U^\perp=\dim V$.
\end{problem}
\begin{solution}
Choose an orthonormal basis of $U$ and extend it to an orthonormal basis of $V$. The added basis vectors form a basis of $U^\perp$, so the dimensions add.
\end{solution}

\begin{problem}[Inner-product synthesis]\label{prob:i13-dossier-12}
Why is an inner product more structure than a vector-space norm?
\end{problem}
\begin{solution}
An inner product determines a norm, but not every norm comes from an inner product. The parallelogram identity characterizes inner-product norms. Thus orthogonality and angle require extra geometric structure beyond linear operations.
\end{solution}

% END VOLUME01 SOLVED DOSSIERS

\section{Exercises with complete solutions}

\begin{exercise}\label{exr:i13-01}
Show nonzero orthogonal vectors are independent.
\end{exercise}
\begin{hint}
Write the relevant inner products explicitly and use bilinearity or sesquilinearity before simplifying norms or angles. Verify positivity, conjugate symmetry, and linearity in the correct argument; do not assume symmetry in the complex case.
\end{hint}
\begin{solution}
Take a dependence relation and inner-product with one vector; all cross terms vanish, leaving its coefficient times a positive norm squared.
\end{solution}

\begin{exercise}\label{exr:i13-02}
Compute the angle between $(1,0)$ and $(1,1)$ in $\mathbb R^2$.
\end{exercise}
\begin{hint}
Compute the norm from \(\|v\|^2=\langle v,v\rangle\) before using an angle formula.
\end{hint}
\begin{solution}
The cosine is $1/(1\cdot\sqrt2)=1/\sqrt2$, so the angle is $\pi/4$.
\end{solution}

\begin{exercise}\label{exr:i13-03}
Find the orthogonal complement of the $x$-axis in $\mathbb R^2$.
\end{exercise}
\begin{hint}
Use the orthogonal decomposition \(v=P_Uv+(v-P_Uv)\) and identify where the residual lies. To show orthogonality, compute the inner product directly and simplify before normalizing anything.
\end{hint}
\begin{solution}
It is the $y$-axis.
\end{solution}

\begin{exercise}\label{exr:i13-04}
Verify the parallelogram identity.
\end{exercise}
\begin{hint}
Find \(U^\perp\) by imposing \(\langle x,u_j\rangle=0\) for a basis \(u_j\) of \(U\).
\end{hint}
\begin{solution}
Expand $\|u+v\|^2+\|u-v\|^2$; the cross terms cancel, leaving $2\|u\|^2+2\|v\|^2$.
\end{solution}

\begin{exercise}\label{exr:i13-05}
Why is complex conjugation necessary in the standard complex inner product?
\end{exercise}
\begin{hint}
Write the relevant inner products explicitly and use bilinearity or sesquilinearity before simplifying norms or angles. For Cauchy--Schwarz, consider \(\|u-tv\|^2\ge0\) and choose \(t\) to minimize the quadratic.
\end{hint}
\begin{solution}
Without conjugation, $\langle v,v\rangle$ need not be real or positive; conjugation gives $\sum|v_i|^2$.
\end{solution}

\begin{exercise}\label{exr:i13-06}
Show $U\cap U^\perp=\{0\}$.
\end{exercise}
\begin{hint}
Use the orthogonal decomposition \(v=P_Uv+(v-P_Uv)\) and identify where the residual lies. Equality in Cauchy--Schwarz occurs exactly when the two vectors are linearly dependent.
\end{hint}
\begin{solution}
If $v$ lies in both, then $\langle v,v\rangle=0$, hence $v=0$ by positive definiteness.
\end{solution}

\begin{exercise}\label{exr:i13-07}
Prove the triangle inequality from Cauchy–Schwarz.
\end{exercise}
\begin{hint}
Write the relevant inner products explicitly and use bilinearity or sesquilinearity before simplifying norms or angles. Use orthogonality to kill the cross term when expanding \(\|u+w\|^2\).
\end{hint}
\begin{solution}
Expand $\|u+v\|^2$ and bound the real part of $\langle u,v\rangle$ by $\|u\|\|v\|$.
\end{solution}

\begin{exercise}\label{exr:i13-08}
Give an example of two different inner products on $\mathbb R^2$.
\end{exercise}
\begin{hint}
Write the relevant inner products explicitly and use bilinearity or sesquilinearity before simplifying norms or angles. For a direct-sum decomposition, prove both orthogonality and that the dimensions add correctly.
\end{hint}
\begin{solution}
The standard product and $\langle x,y\rangle=2x_1y_1+x_2y_2$ are distinct inner products.
\end{solution}


% BEGIN VOL01-EXPANSION I13-exercises-01
\section{Graded supplementary exercises}
These exercises extend the preserved foundational set and are graded by mathematical role.

\subsection*{Standard computations}

\begin{exercise}[Weighted norm]\label{exr:i13-expand-01}
For \(\langle x,y\rangle=2x_1y_1+3x_2y_2\), compute the norm of \((1,2)\).
\end{exercise}
\begin{hint}
Evaluate \(\langle x,x\rangle\) and take the nonnegative square root.
\end{hint}
\begin{solution}
\(\|(1,2)\|^2=2+12=14\), so the norm is \(\sqrt{14}\).
\end{solution}

\begin{exercise}[Angle]\label{exr:i13-expand-02}
Find the angle between \((1,1,0)\) and \((1,0,1)\).
\end{exercise}
\begin{hint}
Compute the dot product and both norms, then use \(\cos\theta=\langle u,v\rangle/(\|u\|\|v\|)\).
\end{hint}
\begin{solution}
The dot product is \(1\), both norms are \(\sqrt2\), so \(\cos\theta=1/2\) and \(\theta=\pi/3\).
\end{solution}

\begin{exercise}[Orthogonal complement]\label{exr:i13-expand-03}
Find \(\operatorname{span}(1,2,2)^\perp\subset\mathbb R^3\).
\end{exercise}
\begin{hint}
Solve \(x+2y+2z=0\) and parametrize the solution plane.
\end{hint}
\begin{solution}
For example \(x=-2s-2t\), so a basis is \((-2,1,0),(-2,0,1)\).
\end{solution}

\begin{exercise}[Orthogonal decomposition]\label{exr:i13-expand-04}
Decompose \(v=(3,1)\) into components parallel and perpendicular to \(u=(1,1)\).
\end{exercise}
\begin{hint}
Project onto the line through \(u\), then subtract.
\end{hint}
\begin{solution}
The parallel component is \((2,2)\); the perpendicular residual is \((1,-1)\).
\end{solution}

\begin{exercise}[Complex product]\label{exr:i13-expand-05}
With \(\langle z,w\rangle=\sum z_i\overline{w_i}\), compute \(\langle(1,i),(i,1)\rangle\).
\end{exercise}
\begin{hint}
Conjugate the entries of the second vector.
\end{hint}
\begin{solution}
\(1(-i)+i(1)=0\), so the vectors are orthogonal.
\end{solution}

\subsection*{Proofs}

\begin{exercise}[Orthogonal independence]\label{exr:i13-expand-06}
Prove every nonzero orthogonal list is linearly independent.
\end{exercise}
\begin{hint}
Take an inner product of a dependence relation with one list vector.
\end{hint}
\begin{solution}
All cross terms vanish, leaving \(a_j\|v_j\|^2=0\); positivity forces every \(a_j=0\).
\end{solution}

\begin{exercise}[Triangle inequality]\label{exr:i13-expand-07}
Derive the triangle inequality from Cauchy--Schwarz.
\end{exercise}
\begin{hint}
Expand \(\|u+v\|^2\) and bound the real part of the cross term.
\end{hint}
\begin{solution}
\(\|u+v\|^2\le\|u\|^2+2\|u\|\|v\|+\|v\|^2=(\|u\|+\|v\|)^2\).
\end{solution}

\begin{exercise}[Double orthogonal complement]\label{exr:i13-expand-08}
For finite-dimensional \(V\), prove \((U^\perp)^\perp=U\).
\end{exercise}
\begin{hint}
Use \(U\subseteq(U^\perp)^\perp\) and compare dimensions.
\end{hint}
\begin{solution}
Both inclusions and \(\dim U+\dim U^\perp=\dim V\) give equal dimensions, hence equality.
\end{solution}

\begin{exercise}[Pythagoras converse]\label{exr:i13-expand-09}
Over a real inner-product space, prove \(\|u+v\|^2=\|u\|^2+\|v\|^2\) iff \(u\perp v\).
\end{exercise}
\begin{hint}
Expand the left side and compare cross terms.
\end{hint}
\begin{solution}
The difference is \(2\langle u,v\rangle\); equality holds exactly when that inner product is zero.
\end{solution}

\subsection*{Counterexamples and hypothesis testing}

\begin{exercise}[Positive weights matter]\label{exr:i13-expand-10}
Show \(\langle x,y\rangle=x_1y_1-x_2y_2\) is not an inner product on \(\mathbb R^2\).
\end{exercise}
\begin{hint}
Test positive definiteness on a coordinate vector.
\end{hint}
\begin{solution}
For \(e_2\), \(\langle e_2,e_2\rangle=-1<0\), so positivity fails.
\end{solution}

\begin{exercise}[Norm without inner product]\label{exr:i13-expand-11}
Show the \(\ell^1\)-norm on \(\mathbb R^2\) does not come from an inner product.
\end{exercise}
\begin{hint}
Test the parallelogram identity on \(e_1,e_2\).
\end{hint}
\begin{solution}
\(\|e_1+e_2\|_1^2+\|e_1-e_2\|_1^2=8\), while \(2\|e_1\|_1^2+2\|e_2\|_1^2=4\).
\end{solution}

\begin{exercise}[Conjugation placement]\label{exr:i13-expand-12}
Explain why \(\sum z_iw_i\) is not the standard complex inner product.
\end{exercise}
\begin{hint}
Evaluate the would-be norm of \((i,0,\ldots)\).
\end{hint}
\begin{solution}
It gives \(i^2=-1\), violating positive definiteness.
\end{solution}

\subsection*{Applications and investigations}

\begin{exercise}[Distance to a plane]\label{exr:i13-expand-13}
Find the distance from \(v=(1,2,3)\) to the plane \(x+y+z=0\).
\end{exercise}
\begin{hint}
The normal vector is \(n=(1,1,1)\); the distance is the magnitude of the projection onto its unit direction.
\end{hint}
\begin{solution}
The distance is \(|1+2+3|/\sqrt3=2\sqrt3\).
\end{solution}

\begin{exercise}[Correlation geometry]\label{exr:i13-expand-14}
For centered data vectors \(x,y\), interpret \(\langle x,y\rangle/(\|x\|\|y\|)\).
\end{exercise}
\begin{hint}
Recognize it as the cosine of the angle between the data vectors.
\end{hint}
\begin{solution}
It is normalized correlation: \(1\) means same direction, \(0\) orthogonality, and \(-1\) opposite direction.
\end{solution}

\subsection*{Challenge problems}

\begin{exercise}[Equality in triangle inequality]\label{exr:i13-expand-15}
Characterize equality in \(\|u+v\|\le\|u\|+\|v\|\).
\end{exercise}
\begin{hint}
Trace equality through Cauchy--Schwarz and the real-part bound.
\end{hint}
\begin{solution}
Equality holds when \(u,v\) point in the same nonnegative scalar direction (with the zero-vector cases included).
\end{solution}

\begin{exercise}[Dimension lower bound]\label{exr:i13-expand-16}
If \(U,W\subset\mathbb R^7\) have dimensions \(5\) and \(4\), prove there is a nonzero vector in \(U\cap W\), then interpret orthogonally.
\end{exercise}
\begin{hint}
Use the dimension formula and \(\dim(U+W)\le7\).
\end{hint}
\begin{solution}
\(\dim(U\cap W)\ge5+4-7=2\). Equivalently, their orthogonal complements have dimensions \(2,3\) and cannot account for all seven directions independently.
\end{solution}
% END VOL01-EXPANSION I13-exercises-01
\section*{Chapter summary}
Inner products add geometry to linear algebra. Length, angle, orthogonality and best approximation all come from a positive definite sesquilinear form. The complex case forces careful placement of conjugation, but the resulting theory unifies Euclidean geometry with finite-dimensional Hilbert-space methods.
The important habit is to move fluently between invariant statements and concrete coordinates while remembering which conclusions depend on finite dimensionality, the scalar field, or the inner product.
